\chapter{Data Analysis with Pandas and Visualisation with Matplotlib}
\label{ch:pandas_matplotlib}

\section{Introduction to Tabular Data Analysis with Pandas}

Pandas is the premier Python library for structured data manipulation. Built atop NumPy, Pandas introduces labeled, heterogeneous tabular structures that support relational database operations, time series analysis, and statistical processing.

\subsection{Primary Data Structures: Series and DataFrame}
\begin{enumerate}[leftmargin=1.8em, itemsep=3pt]
    \item \textbf{Series:} A 1-dimensional labeled homogeneous array capable of holding any data type. It consists of an underlying data array and an associated \textbf{Index}.
    \item \textbf{DataFrame:} A 2-dimensional tabular structure with labeled axes (rows and columns). It can be viewed as an ordered dictionary of \texttt{Series} sharing a common index.
\end{enumerate}

\begin{lstlisting}[caption={Creating and Manipulating DataFrames}]
import pandas as pd
import numpy as np

# Constructing DataFrame from dictionary of lists
data = {
    'Timestamp': pd.date_range(start='2026-01-01', periods=5, freq='h'),
    'Node_ID': ['N1', 'N2', 'N1', 'N3', 'N2'],
    'Voltage_V': [3.31, 4.98, 3.28, np.nan, 5.02],
    'Current_mA': [12.4, 45.1, 11.9, 8.5, 46.0]
}
df = pd.DataFrame(data)
print(df.info())
\end{lstlisting}

\subsection{Selection and Indexing: \texttt{.loc[]} vs. \texttt{.iloc[]}}
\begin{itemize}[leftmargin=1.8em, itemsep=3pt]
    \item \textbf{\texttt{.loc[]}: Label-based indexing.} Slicing with \texttt{.loc} includes \textbf{both} the start and the stop endpoints!
    \item \textbf{\texttt{.iloc[]}: Integer-position-based indexing.} Slicing follows standard Python sequence semantics (half-open interval $[start, stop)$, excluding the stop boundary).
\end{itemize}

\begin{lstlisting}[caption={Label vs Position Slicing}]
# Select rows 0 to 2 and specific columns by label
sub_loc = df.loc[0:2, ['Node_ID', 'Voltage_V']]

# Select rows 0 to 2 (indices 0 and 1) and columns 1 to 3 by position
sub_iloc = df.iloc[0:2, 1:3]
\end{lstlisting}

\subsection{Data Cleaning: Handling Missing Values}
Real-world engineering sensor streams contain missing data resulting from packet drops or sensor failures:
\begin{lstlisting}[caption={Handling Missing Data}]
# Check missing counts
print(df.isna().sum())

# Drop rows with any NaN
df_clean = df.dropna()

# Impute missing values with group mean
df['Voltage_V'] = df['Voltage_V'].fillna(df['Voltage_V'].mean())
\end{lstlisting}

\subsection{Aggregation and Grouping: Split-Apply-Combine}
\begin{lstlisting}[caption={Group-by Aggregations}]
summary = df.groupby('Node_ID').agg({
    'Voltage_V': ['mean', 'std'],
    'Current_mA': ['mean', 'max']
})
print(summary)
\end{lstlisting}

\section{Engineering Data Visualisation with Matplotlib}

Matplotlib is a multi-platform data visualization library built on NumPy arrays. It provides full control over every graphic primitive on the canvas.

\subsection{Matplotlib Object-Oriented Hierarchy}
While the \texttt{pyplot} interface mimics MATLAB procedural commands, the object-oriented API provides robust control:
\begin{itemize}[leftmargin=1.8em, itemsep=2pt]
    \item \textbf{Figure:} The top-level container holding all canvas elements, subplots, titles, and legends.
    \item \textbf{Axes:} The actual plotting region bounding the data coordinate system with $x/y$ scales, ticks, and labels.
\end{itemize}

\begin{lstlisting}[caption={Engineering Signal Visualization using the OO API}]
import matplotlib.pyplot as plt
import numpy as np

# Synthetic sensor telemetry
t = np.linspace(0.0, 2.0, 500)
freq = 5.0
signal_clean = 3.3 * np.sin(2 * np.pi * freq * t)
noise = 0.4 * np.random.randn(len(t))
signal_noisy = signal_clean + noise

# Create Figure and Axes
fig, (ax1, ax2) = plt.subplots(2, 1, figsize=(10, 6), sharex=True)

# Subplot 1: Time domain waveform
ax1.plot(t, signal_noisy, color='gray', alpha=0.6, label='Noisy ADC Signal')
ax1.plot(t, signal_clean, color='crimson', lw=2, label='True Fundamental (5 Hz)')
ax1.set_ylabel('Voltage (V)', fontsize=11)
ax1.set_title('Sensor Telemetry Acquisition: Time Domain Waveform', fontsize=12, fontweight='bold')
ax1.grid(True, linestyle='--', alpha=0.7)
ax1.legend(loc='upper right')

# Subplot 2: Error residual distribution
residual = signal_noisy - signal_clean
ax2.plot(t, residual, color='teal', lw=1.2, label='Noise Residual')
ax2.axhline(0, color='black', linestyle=':', lw=1)
ax2.set_xlabel('Time (s)', fontsize=11)
ax2.set_ylabel('Residual (V)', fontsize=11)
ax2.grid(True, linestyle='--', alpha=0.7)
ax2.legend(loc='upper right')

plt.tight_layout()
# Save high-resolution vector graphic
# plt.savefig('sensor_telemetry.pdf', dpi=300)
\end{lstlisting}

\section{Standard Statistical Plot Types}
\begin{enumerate}[leftmargin=1.8em, itemsep=3pt]
    \item \textbf{Line Plot (\texttt{ax.plot}):} Displays continuous time-series trajectories.
    \item \textbf{Scatter Plot (\texttt{ax.scatter}):} Visualizes bi-variate correlation and clustering.
    \item \textbf{Histogram (\texttt{ax.hist}):} Approximates continuous probability density functions.
    \item \textbf{Bar Chart (\texttt{ax.bar}):} Compares discrete categorical quantities.
    \item \textbf{Box Plot (\texttt{ax.boxplot}):} Visualizes 5-number summary (minimum, Q1, median, Q3, maximum) and outlier anomalies.
\end{enumerate}

\section{Chapter Summary and Review Questions}
\begin{enumerate}[leftmargin=1.8em, itemsep=4pt]
    \item Compare and contrast the endpoint indexing behavior of \texttt{.loc[]} versus \texttt{.iloc[]}.
    \item Explain the three stages of the Split-Apply-Combine strategy implemented by Pandas \texttt{groupby()}.
    \item Why is the Object-Oriented API (\texttt{fig, ax = plt.subplots()}) preferred over procedural state-machine calls (\texttt{plt.plot()}) in professional engineering workflows?
    \item Given a DataFrame \texttt{df} with missing values in column \texttt{'temperature'}, write code to replace missing entries with the rolling median of the preceding 3 observations.
\end{enumerate}
